\section*{第 \ref{ch:g12:nuclear-energy} 章——核能：裂变、聚变与 $E=mc^2$}

\begin{solution}{exo:g12:nuclear-energy:1}
(a) $\qty{1}{eV} = \qty{1.60e-19}{J}$;
$\qty{1}{MeV} = \qty{1.60e-13}{J}$.
(b) $200 \times \num{1.60e-13} = \qty{3.2e-11}{J}$.
(c) $E = \num{1.66054e-27} \times \num{9.00e16} = \qty{1.49e-10}{J}
= \num{1.49e-10}/\num{1.60e-13} \approx \qty{9.3e2}{MeV}$——与课本中的
\qty{931.5}{MeV} 相符，在所取 $c$ 的近似精度之内。
\end{solution}

\begin{solution}{exo:g12:nuclear-energy:2}
$E = \num{6.0e-3} \times \num{9.00e16} = \qty{5.4e14}{J}$;
$\num{5.4e14}/\num{3.6e7} = \num{1.5e7}$ 天 $\approx \num{4.1e4}$
年——一块方糖的能量，可供四万年之用。
\end{solution}

\begin{solution}{exo:g12:nuclear-energy:3}
$\Delta m = 6 \times 1.00728 + 6 \times 1.00866 - 11.99671 =
\qty{0.09893}{u}$; $E_b = 0.09893 \times 931.5 \approx \qty{92.2}{MeV}$;
$E_b/A \approx \qty{7.68}{MeV}$ 每个核子。
\end{solution}

\begin{solution}{exo:g12:nuclear-energy:4}
(a) $A = 236 - 144 - 3 = 89$, $Z = 92 - 56 = 36$:
${}^{89}_{36}\mathrm{Kr}$，氪。
(b) $A = 236 - 97 - 2 = 137$, $Z = 92 - 40 = 52$:
${}^{137}_{52}\mathrm{Te}$，碲。
\end{solution}

\begin{solution}{exo:g12:nuclear-energy:5}
比结合能约为 $7.1$、$7.7$、$8.8$ 和 \qty{7.6}{MeV}；铁-56 是
结合得最紧密的。铀可通过裂变放出能量，氦-4 与
碳-12 则通过聚变放出能量——两种过程都向铁靠拢；铁本身
已无处可去。
\end{solution}

\begin{solution}{exo:g12:nuclear-energy:6}
$\Delta m = 1.00728 + 1.00866 - 2.01355 = \qty{0.00239}{u}$，因此
$E_b \approx \qty{2.23}{MeV}$——约为化学键中几
\unit{eV} 的一百万倍。
\end{solution}

\begin{solution}{exo:g12:nuclear-energy:7}
之前：$\qty{236.00216}{u}$；之后：$143.89220 + 88.89790 +
3 \times 1.00866 = \qty{235.81608}{u}$; $\Delta m = \qty{0.18608}{u}$,
因此 $E \approx \qty{173}{MeV} = 173 \times \num{1.60e-13} \approx
\qty{2.8e-11}{J}$.
\end{solution}

\begin{solution}{exo:g12:nuclear-energy:8}
$\Delta m = 2 \times 2.01355 - 3.01493 - 1.00866 = \qty{0.00351}{u}$:
$E \approx \qty{3.27}{MeV}$。D--T 反应生成氦-4，其结合异常紧密
（每个核子 \qty{7.1}{MeV}）；氦-3 在曲线上位置低得多，因此
该步释放的能量少得多。
\end{solution}

\begin{solution}{exo:g12:nuclear-energy:9}
(a) 中子不带电：没有库仑势垒，裂变在低温下也能发生。
聚变的原子核都带正电，必须翻越势垒，因此只有
$\sim \qty{1e8}{K}$ 的热运动才能使碰撞足够剧烈。(b) 核心的质子数是天文数字；
其中极少数最快的质子负责聚变，
而太阳有数十亿年的富余时间，
恒星需要的正是缓慢燃烧。(c) 不存在中子雪崩：
每次聚变都要靠自身的碰撞来换取，任何约束的丧失
都会使等离子体降温，火便熄灭。
\end{solution}

\begin{solution}{exo:g12:nuclear-energy:10}
$\Delta m = 4 \times 1.00728 - 4.00151 = \qty{0.02761}{u}$，因此
$E \approx \qty{25.7}{MeV}$；转化的质量占比为
$0.02761/4.02912 \approx 0.7\%$。
\end{solution}

\begin{solution}{exo:g12:nuclear-energy:11}
(a) 它通过碰撞使裂变中子减速；慢中子
使铀-235 裂变的可能性大得多，因此链式反应需要它们。(b) 控制棒
吸收中子，使链式反应维持在每次裂变仅余一个中子的状态——
插入过深则反应熄灭。(c) 无水、无慢化：中子
保持高速，难以被捕获，链式反应窒息——设计失效时
趋向“关闭”。
\end{solution}

\begin{solution}{exo:g12:nuclear-energy:12}
每代裂变数翻倍，因此经过 $n$ 代后，约有 $2^n$
个原子核发生了裂变；$2^{n} = \num{2.56e21} \approx 2 \times (10^{3})^{7}
\approx 2^{71}$ 得 $n \approx 71$。时间：$71 \times \qty{10}{ns}
\approx \qty{0.7}{\micro s}$——一克燃料，相当于三吨煤的
能量，在不到一微秒内释放：临界是悬崖，而非斜坡。
\end{solution}

\begin{solution}{exo:g12:nuclear-energy:13}
$\Delta m = 55.92066 - 2 \times 27.96924 = \qty{-0.01782}{u}$:
$E \approx \qty{-16.6}{MeV}$——裂变\emph{吸收}能量。
铁位于峰值：使其裂变或聚变，结合能都会下降，
两者都需付出能量。它是一切核火的灰烬。
\end{solution}

\begin{solution}{exo:g12:nuclear-energy:14}
Fe: $\Delta m = 26 \times 1.00728 + 30 \times 1.00866 - 55.92066 =
\qty{0.52842}{u}$, $E_b \approx \qty{492}{MeV}$, $E_b/A \approx
\qty{8.79}{MeV}$. U: $\Delta m = 92 \times 1.00728 + 143 \times
1.00866 - 234.99350 = \qty{1.91464}{u}$, $E_b \approx \qty{1784}{MeV}$,
$E_b/A \approx \qty{7.59}{MeV}$。将 $235$ 个核子重新排列，
若按每个核子 \qty{8.5}{MeV} 而非 \qty{7.59}{MeV} 计，可增加 $235 \times 0.9
\approx \qty{2e2}{MeV}$——与例题中的 $185$--$200$ 吻合。
\end{solution}

\begin{solution}{exo:g12:nuclear-energy:15}
$\num{5.0e8}/\num{2.82e-12} \approx \num{1.8e20}$ 次反应，每
秒；质量变化率 $\num{1.8e20} \times 5.029 \times \num{1.66054e-27}
\approx \qty{1.5e-6}{kg/s}$，即每天约 \qty{0.13}{kg}——
与 \num{1.4e6} 千克煤相比，相差一千万倍。
\end{solution}

\begin{solution}{pb:g12:nuclear-energy:1}
\textbf{1.} $P_{\text{热}} = \num{1.0e9}/0.33 \approx \qty{3.0}{GW}$。

\textbf{2.} $200 \times \num{1.60e-13} = \qty{3.2e-11}{J}$.

\textbf{3.} $\num{3.0e9}/\num{3.2e-11} \approx \num{9.4e19}$ 次裂变
每秒。

\textbf{4.} $m_{\mathrm U} = 235 \times \num{1.66054e-27} =
\qty{3.90e-25}{kg}$; $\num{9.4e19} \times \num{3.90e-25} \approx
\qty{3.7e-5}{kg/s}$——每秒约四十微克。

\textbf{5.} $\num{3.7e-5} \times \num{3.16e7} \approx \qty{1.2e3}{kg}$:
每年约 $1.2$ 吨铀-235。

\textbf{6.} $\num{1.2e3}/0.040 \approx \qty{2.9e4}{kg}$——约 $29$
吨燃料：一辆重型卡车，每年一趟。

\textbf{7.} $\num{3.0e9} \times \num{3.16e7} \approx \qty{9.5e16}{J}$.

\textbf{8.} $\num{9.5e16}/\num{3.0e7} \approx \qty{3.2e9}{kg}$: $3.2$
百万吨煤。

\textbf{9.} $\num{3.2e9}/\num{1.2e3} \approx \num{2.7e6}$——恰为
$\num{8.2e13}/\num{3.0e7}$ 的阶梯之比，理应如此。

\textbf{10.} $\num{3.2e9}/\num{3.0e6} \approx 1060$ 列火车，每年
约每天三列：一列运煤列车于黎明、正午和黄昏发车，永不间断，
相比之下，每年春季只需一辆卡车。

\textbf{11.} $17.6 \times \num{1.60e-13} = \qty{2.82e-12}{J}$.

\textbf{12.} 一对核的质量为 $5.029 \times \num{1.66054e-27} =
\qty{8.35e-27}{kg}$: $\num{2.82e-12}/\num{8.35e-27} \approx
\qty{3.4e14}{J/kg}$——即阶梯上的聚变梯级。

\textbf{13.} $\num{9.5e16}/\num{3.4e14} \approx \qty{2.8e2}{kg}$ 的 D--T
燃料；其中氘的份额为 $280 \times 2.014/5.029 \approx \qty{1.1e2}{kg}$。

\textbf{14.} $\num{1.1e2}/\num{3.3e-2} \approx \qty{3.4e3}{m^3}$ 的海水
——约一个半游泳池，供全市一年之用。

\textbf{15.} 不会自我倍增：没有中子点燃下一次聚变，
一旦约束失效，等离子体便会降温，并在数秒内停止。
难点恰恰相反——要把 \qty{1.5e8}{K} 维持住，
使其持续足够长的时间以便燃烧。

\textbf{16.} $\num{3.9e26}/\num{9.00e16} \approx \qty{4.3e9}{kg/s}$:
每秒四百万吨。

\textbf{17.} $\num{3.0e9}/\num{9.00e16} \approx \qty{3.3e-8}{kg/s}$,
即 $\num{3.3e-8} \times \num{3.16e7} \approx \qty{1.0}{kg}$ 每年
——确实如此：$\qty{1.2e3}{kg} \times 0.00085 \approx 1$，$\num{3.2e9}
\times \num{3.3e-10} \approx 1$，$280 \times 0.00375 \approx 1$。每种
燃料都释放同样的一千克；它们的差别只在于
需要多少货物来运送它。

\textbf{18.} $\num{4.3e9} \times \num{3.16e7} \times \num{4.5e9}
\approx \qty{6.1e26}{kg}$——仅为 $\num{6.1e26}/\num{2.0e30} \approx
\num{3e-4}$，即万分之三。

\textbf{19.} 可转化的质量 $0.0007 \times \num{2.0e30} =
\qty{1.4e27}{kg}$，相当于 $\num{1.4e27} \times \num{9.00e16} \approx
\qty{1.3e44}{J}$；在 \qty{3.9e26}{W} 的功率下，可维持 $\num{1.3e44}/
\num{3.9e26} \approx \qty{3.2e17}{s} \approx \num{1.0e10}$ 年——
一百亿年，其中约五十亿年尚余。

\textbf{20.} 全市一年的用量：$3.2$ 百万吨煤，或
$\num{1.2e3}$ 千克铀-235，或 $\num{2.8e2}$ 千克
D--T，取自约一个半游泳池的海水。在每一种账目中，
同样的一笔：一千克质量，已消失。$E = mc^2$ 并不在意
究竟由哪种燃料携带它——只在乎需要多大的卡车。
\end{solution}
